\section{Strictly compatible polynomials of the raw core}
\label{sec:compatible}

Let
\[
\mathsf M_0=(X_{4,0},\pi_{\core,0},0).
\]
Spread \(X_{4,0}\), the hyperplane class, and the rational correspondence
\(\pi_{\core,0}\) over the complement of a finite set \(S\) of rational
primes.  At \(p\notin S\), specialization of the projector commutes with
geometric Frobenius \(F_p\).

\subsection{Rationality and coefficient-field independence}

Choose an integer \(D>0\) that clears the denominator of the projector.
For every \(m\ge1\), the trace of
\(\pi_{\core,0}F_p^m\) is a rational multiple of the trace of an integral
algebraic correspondence.  The correspondence comparison theorem of
Katz--Messing makes these traces independent of the chosen Weil
cohomology, hence independent of \(\ell\)
\cite[Theorem 2(2), pp. 76--77]{KatzMessing1974}.  Applying Newton
identities to the first ten traces gives one polynomial
\begin{equation}
\chi_{p,\core}(U)
=\det(U-F_p\mid H_\ell(\mathsf M_0))\in\Q[U]
\label{eq:chi-core-Q}
\end{equation}
of degree \(10\), independent of \(\ell\).

No semisimplicity is used.  The idempotent itself decomposes the
realization into its image and kernel, both stable under Frobenius.

\subsection{The separate integrality step}

Denominator clearing in the cycle does not by itself prove integral
coefficients.  Instead, consider the full middle characteristic polynomial
\[
\chi_{p,H^5}(U)=\det(U-F_p\mid H^5_{\mathrm{\acute et}})
\in\Z[U].
\]
The projector decomposition gives a factorization
\[
\chi_{p,H^5}(U)=
\chi_{p,\core}(U)\chi_{p,\lev}(U)
\]
into monic polynomials in \(\Q[U]\).  Every root of either factor is a
root of the full polynomial and is therefore an algebraic integer.
The rational coefficients of each factor are algebraic integers, hence
integers.  Thus
\[
\chi_{p,\core}(U)\in\Z[U].
\]
Reversing coefficients now gives
\begin{equation}
P_p(T)=\det(1-F_pT\mid H_\ell(\mathsf M_0))
=T^{10}\chi_{p,\core}(T^{-1})\in\Z[T].
\label{eq:P-from-chi}
\end{equation}
The polynomial \(\chi_{p,\core}\) is monic; \(P_p\), whose constant term
is \(1\), is not called monic.

\subsection{Purity and reciprocity}

The roots of \(\chi_{p,\core}\) are eigenvalues on \(H^5\), so they are
pure of weight \(5\).  The middle Poincar\'e pairing takes values in
\(\Q_\ell(-5)\), hence
\[
\langle F_px,F_py\rangle=p^5\langle x,y\rangle.
\]
Self-transposition and idempotence give
\[
\ker\pi_{\core,0}=(\operatorname{im}\pi_{\core,0})^\perp,
\]
so the restricted pairing is nondegenerate.  The spread projector commutes
with \(F_p\), and its image is stable.  Frobenius is therefore a similitude
of multiplier \(p^5\) on the core, and its ten eigenvalues occur in pairs
\[
\alpha,\quad \frac{p^5}{\alpha}.
\]
Their product is \(p^{25}\), and
\begin{equation}
P_p(T)=p^{25}T^{10}P_p\left(\frac1{p^5T}\right).
\label{eq:reciprocity}
\end{equation}
Writing \(P_p(T)=\sum_{k=0}^{10}a_kT^k\) and comparing coefficients yields
\[
a_{10-k}=p^{25-5k}a_k.
\]
This argument uses the invariant direct summand supplied by the idempotent,
not a semisimplicity hypothesis.

Finally, geometric Frobenius acts as \(p\) on \(\Q_\ell(-1)\), hence as
\(p^{-2}\) on \(\Q_\ell(2)\).  The normalized core
\(\mathsf M_0(2)\) has weight \(1\) and polynomial
\[
P_p^{(2)}(T)=P_p(T/p^2)\in\Q[T].
\]
Its denominators are the prescribed Tate denominators, not a failure of
compatibility.  This completes Theorem~\ref{thm:compatible-main}.
